\documentclass{beamer}

\mode<presentation>





\usetheme{Frankfurt}%
\usecolortheme{seagull}
\logo{\includegraphics[height=.25in]{clarksonGreen}}

\definecolor{garnet}{RGB}{136,0,0}
%\definecolor{clarksonGreen}{RGB}{0,71,28}
\definecolor{clarksonGreen}{RGB}{0,52,21}
\setbeamercolor{palette primary}{fg=clarksonGreen,bg=white}
\setbeamercolor{palette secondary}{fg=clarksonGreen,bg=white}
\setbeamercolor{palette tertiary}{fg=clarksonGreen,bg=white}
\setbeamercolor{palette quaternary}{bg=clarksonGreen,fg=white}
\setbeamercolor{block title}{fg=black,bg=black!15}
\setbeamercolor{block body}{fg=black,bg=black!10}
\setbeamercolor{titlelike}{bg=clarksonGreen,fg=white} % parent=palette quaternary}

\newcommand{\half}{\mbox{$\frac{1}{2}$}}
\newcommand{\deltat}{\mbox{$\triangle t$}}
\newcommand{\deltax}{\mbox{$\triangle x$}}
\newcommand{\deltay}{\mbox{$\triangle y$}}

\newcommand{\deriv}[2]{\frac{d}{d#2}#1}
\newcommand{\derivTwo}[2]{\frac{d^2}{d#2^2}#1}

\newcommand{\lp}{\left(}
\newcommand{\rp}{\right)}

\newcommand{\laplace}[1]{\makebox{$ {\cal L} \left\{ #1 \right\}$}}
\newcommand{\invlaplace}[1]{\makebox{$ {\cal L}^{-1} \left\{ #1 \right\}$}}

\begin{document}

\title{Ordinary Differential Equations}
\subtitle{Math 232 - Week 13, Day 1}

\author{Kelly Black}
\institute{Clarkson University}
\date{21 November 2011}

\begin{frame}
  \titlepage
\end{frame}

\begin{frame}
  \frametitle{Outline}
  \tableofcontents[pausesection,hideallsubsections]
\end{frame}


\section{Laplace Transform Table}

\begin{frame}
  \frametitle{Backwards Going Laplace Transforms }

  Before we begin a quick recap:
  \begin{enumerate}
  \item $\laplace{c_1 f(t) + c_2 g(t)}=c_1 \laplace{f(t)} + c_2 \laplace{g(t)}$
  \item $\laplace{f}=F(s)$
  \item $\laplace{t^nf(t)}=(-1)^n\frac{d^n}{ds^n}F(s)$
  \end{enumerate}

  Also:
  \begin{eqnarray*}
    \laplace{f(t)}  & = & \int^\infty_0 f(t) e^{-st} ~ dt \\ 
    \laplace{t^n}    & = & \frac{n!}{s^{n+1}} \\ 
    \laplace{e^{at}} & = & \frac{1}{s-a} \\ 
    \laplace{\sin(\omega t)} & = & \frac{\omega}{s^2+\omega^2} \\
    \laplace{\cos(\omega t)} & = & \frac{s}{s^2+\omega^2} \\ 
    \laplace{e^{at}f(t)} & = & \laplace{f}(s-a) 
  \end{eqnarray*}


\end{frame}

\section{The Inverse Laplace Transform}


\begin{frame}
  \frametitle{The Inverse Laplace Transform}

  \begin{block}{Question:}
    If $\laplace{f} = \frac{1}{s+3}$ what is f(t)?
  \end{block}

  \uncover<2->
  {
    Well, $a=-3$ so
    \begin{eqnarray*}
      \laplace{e^{at}} & = & \frac{1}{s-a}, \\
      f(t) & = & e^{-3t}.
    \end{eqnarray*}
  }

  \uncover<3->
  {
    Notation:
    \begin{eqnarray*}
      \invlaplace{\frac{1}{s+3}} & = & e^{-3t}.
    \end{eqnarray*}
  }

\end{frame}


\begin{frame}
  \frametitle{Example}

  \begin{block}{Question:}
    If $\laplace{f} = \frac{4}{s+3}-\frac{2}{s-5}$ what is f(t)?
  \end{block}

  \uncover<2->
  {
    \begin{eqnarray*}
      f & = & 4 e^{-3t} - 2 e^{5t}.
    \end{eqnarray*}
  }

\end{frame}


\begin{frame}
  \frametitle{Same Example, Only Different}

  \begin{block}{Question:}
    If $\laplace{f} = \frac{2s-26}{s^2-2s-15}$ what is f(t)?
  \end{block}

  \uncover<2->
  {
    \begin{eqnarray*}
      \frac{2s-26}{s^2-2s-15} & = & \frac{2s-26}{(s+3)(s-5)}, \\
      & = & \frac{a}{s+3} + \frac{b}{s-5}, \\
      & = & \frac{4}{s+3} - \frac{2}{s-5}.
    \end{eqnarray*}
  }

  \uncover<3->
  {
    \begin{eqnarray*}
      f & = & 4 e^{-3t} - 2 e^{5t}.
    \end{eqnarray*}
  }


\end{frame}


\begin{frame}
  \frametitle{Example}

  \begin{block}{Question:}
    If $\laplace{f} = \frac{-2}{(s-2)^4}$ what is f(t)?
  \end{block}

  \uncover<2->
  {
    \begin{eqnarray*}
      \laplace{e^{at}f(t)} & = & F(s-a), \\
      \Rightarrow g(t) & = & e^{2t} f(t), \\
      \laplace{f} & = & \frac{-2}{s^4}
    \end{eqnarray*}
  }



\end{frame}


\begin{frame}

  
    \begin{eqnarray*}
      \laplace{t^n} & = & \frac{n!}{s^{n+1}}, \\ [10pt]
      \Rightarrow n & = & 3, \\
      \laplace{f} & = & \frac{-2}{s^{4}}, \\
      & = & \frac{-2}{3!}\cdot\frac{3!}{s^4}, \\ [10pt]
      f(t) & = & \frac{-2}{3!} t^3, \\ [10pt]
      g(t) & = & -\frac{1}{3} e^{2t} t^3.
    \end{eqnarray*}


\end{frame}


\begin{frame}
  \frametitle{Example}

  \begin{block}{Question:}
    If $\laplace{f} = \frac{s+1}{(s+1)^2+9}$ what is f(t)?
  \end{block}

  \uncover<2->
  {
    \begin{eqnarray*}
      \laplace{e^{t}f(t)} & = & \frac{s+1}{(s+1)^2+9}, \\
      \Rightarrow g(t) & = & e^{t} f(t), \\
      \laplace{f} & = & \frac{s}{s^2+9}
    \end{eqnarray*}
  }



\end{frame}


\begin{frame}

    \begin{eqnarray*}
      \laplace{\cos(at)} & = & \frac{s}{s^2+a^2}, \\ [10pt]
      \Rightarrow a & = & 3, \\
      \laplace{f} & = & \frac{s}{s^2+9}, \\
      f(t) & = & \cos(3t), \\ [10pt]
      g(t) & = & e^{t}\cos(3t).
    \end{eqnarray*}

\end{frame}


\begin{frame}
  \frametitle{Example}

  \begin{block}{Question}
    if $\laplace{g}=\frac{s}{s^2+2s+17}$ what is $g$?
  \end{block}

    \uncover<2->
    {
      \begin{eqnarray*}
        \laplace{g} & = & \frac{s}{s^2+2s+17}, \\
        & = & \frac{s}{(s+1)^2+16}, \\
        & = & \frac{s+1}{(s+1)^2+16} - \frac{1}{4}\cdot\frac{4}{(s+1)^2+16}.
      \end{eqnarray*}
    }

    \uncover<3->
    {
      \begin{eqnarray*}
        g(t) & = & e^t\cos(4t) - \frac{1}{4}\sin(4t).
      \end{eqnarray*}
    }
  

\end{frame}


\begin{frame}
  \frametitle{Example}

  \begin{block}{Question}
    if $\laplace{g}=\frac{s^2-14s-58}{s^3+2s^2+21s-58}$ what is $g$?
  \end{block}

    \uncover<2->
    {
      \begin{eqnarray*}
        \laplace{g} & = & \frac{s^2-14s-58}{s^3+2s^2+21s-58}, \\
        & = & \frac{s^2-14s-58}{(s-2)(s^2+4s+29)}, \\
        & = & \frac{as+b}{s^2+4s+29}+\frac{c}{s-2}, \\
        & = & 3\frac{s+2}{(s+2)^2+25} - 3\frac{2}{(s+1)^2+25}-\frac{2}{s-2}.
      \end{eqnarray*}
    }

    \uncover<3->
    {
      \begin{eqnarray*}
        g(t) & = & 3e^{-2t}\cos(5t) - \frac{6}{5}e^{-2t}\sin(5t)-2e^{2t}.
      \end{eqnarray*}
    }
  

\end{frame}


\begin{frame}
  \frametitle{Example}

  \begin{block}{Question}
    if $\laplace{g}=\frac{2s+1}{s^2(s^2+9)}$ what is $g$?
  \end{block}

    \uncover<2->
    {
      \begin{eqnarray*}
        \laplace{g} & = & \frac{2s+1}{s^2(s^2+9)}, \\
        & = & \frac{a}{s}+\frac{b}{s^2} + \frac{cs+d}{s^2+9}, \\
        & = & \frac{2}{9}\cdot\frac{1}{s} + \frac{1}{9}\cdot\frac{1}{s^2}-\frac{2}{9}\cdot\frac{s}{s^2+9}-\frac{1}{27}\cdot\frac{3}{s^2+9}.
      \end{eqnarray*}
    }

    \uncover<3->
    {
      \begin{eqnarray*}
        g(t) & = & \frac{2}{9} + \frac{1}{9}t - \frac{2}{9}\cos(3t) - \frac{1}{27} \sin(3t).
      \end{eqnarray*}
    }
  

\end{frame}


\begin{frame}
  \frametitle{Example}

  \begin{block}{Question}
    if $\laplace{g}=\frac{8s}{(s^2+16)^2}$ what is $g$?
  \end{block}

    \uncover<2->
    {
      \begin{eqnarray*}
        \laplace{g} & = & \frac{8s}{(s^2+16)^2}, \\
        & = & -4 \frac{-2s}{(s^2+16)^2}, \\
        & = & - \frac{d}{ds} \frac{4}{s^2+16}.
      \end{eqnarray*}
    }

    \uncover<3->
    {
      \begin{eqnarray*}
        g(t) & = & t \sin(4t).
      \end{eqnarray*}
    }
  

\end{frame}

\begin{frame}
  \frametitle{Example}

  \begin{block}{Question}
    if $\laplace{g}=\frac{8s^3-5s^2-12}{s^4-16}$ what is $g$?
  \end{block}

    \uncover<2->
    {
      \begin{eqnarray*}
        \laplace{g} & = & \frac{8s^3-5s^2-12}{s^4-16}, \\
        & = & \frac{82^3-5s^2-12}{(s-2)(s+2)(s^2+4)}, \\
        & = & \frac{a}{s-2} + \frac{b}{s+2} + \frac{cs+d}{s^2+4}, \\
        & = & \frac{1}{s-2} + \frac{3}{s+2} + 4 \frac{s}{s^2+4} - \frac{1}{2}\cdot\frac{2}{s^2+4}.
      \end{eqnarray*}
    }

    \uncover<3->
    {
      \begin{eqnarray*}
        g(t) & = & e^{2t} + 3e^{-2t} + 4\cos(2t) - \half\sin(2t).
      \end{eqnarray*}
    }
  

\end{frame}




\end{document}

% LocalWords:  Clarkson pausesection hideallsubsections
